\documentclass{article}
\usepackage{graphicx}
\usepackage{amsmath}
\usepackage{amssymb}
\usepackage{amsfonts}
\usepackage{hyperref}
\usepackage{exercises}
\usepackage{indentfirst}
\usepackage{amsthm}
\usepackage{csquotes}
\usepackage{yhmath}
\usepackage{comment}
\usepackage{placeins}
\usepackage{tikz}
\newcommand*{\oldneg}{\mathord{\sim}}
\newcommand{\Z}{\mathbb{Z}}
\newcommand{\N}{\mathbb{N}}
\let\oldemptyset\emptyset

\theoremstyle{definition}
\newtheorem{definition}{Definition}

\setlength{\parindent}{.5in}
\setlength{\parskip}{12pt}
\title{Chapter 6 Answers}
\begin{document}
\date{}
\maketitle

\section*{Problem 1}
\textit{Note}: Here is why the diagonal function is implied by the iteration theorem. 
Consider any $R_y(x,y,z)$. By the iteration theorem, there is some $\phi(y,z)$ such that 
$R_y(x,y,z)$ iff $R_{\phi(y,z)}(x)$. By definition of $R^n_i(x_1,...,x_n)$ when $n=1$, $R_{\phi(y,z)}(x)$ means $x\in\omega_{\phi(y,z)}$, and thus 
$x\in\omega_{\phi(y,z)}$ iff $R_y(x,y,z)$.


Consider any two recursively enumerable relations $M_1(x,y,z)$ and $M_2(x,y,z)$. Hence the relations $M_1(x,d(y,z),d(z,y))$ and $M_2(x,d(z,y),d(y,z))$ are 
recursively enumerable, and have indices $m$ and $n$, respectively. Thus, if $d(y,z)$ is the diagonal function:
$x\in\omega_{d(m,n)}$ iff $M_1(x,d(m,n),d(n,m))$ and $x\in\omega_{d(n,m)}$ iff $M_2(x,d(m,n),d(n,m)$, thus showing that $a=d(m,n)$ and $b=d(n,m)$ satisfies the 
theorem.

\section*{Problem 2}
Consider any two recursively enumerable relations $M_1(x,y_1,y_2,z_1,z_2)$ and $M_2(x,y_1,y_2,z_1,z_2)$. Hence the relations 
$M_1(x,i,j,d(y,z),d(z,y))$ and $M_2(x,i,j,d(y,z),d(z,y))$ (as relations between $x,y,z,i,j$) are recursively enumerable for the diagonal function $d$. Hence, by the 
iteration theorem, there are recursive functions $t_1(i,j)$ and $t_2(i,j)$ such that for all numbers $i,j,x,y,z$:\newline
(1): $R_{t_1(i,j)}(x,y,z)$ iff $M_1(x,i,j,d(y,z),d(z,y))$;\newline
(2): $R_{t_2(i,j)}(x,y,z)$ iff $M_2(x,i,j,d(z,y),d(y,z))$.\newline
Let $\phi_1(i,j)$ be the function $d(t_1(i,j),t_2(i,j))$ and $\phi_2(i,j)$ be the function $d(t_2(i,j),t_1(i,j))$. Thus
$x\in\omega_{\phi_1(i,j)}$ iff $R_{t_1(i,j)}(x,t_1(i,j),t_2(i,j))$ iff $M_1(x,i,j,d(t_1(i,j),t_2(i,j)),d(t_2(i,j),t_1(i,j))$ iff 
$M_1(x,i,j,\phi_1(i,j),\phi_2(i,j))$ and 
$x\in\omega_{\phi_2(i,j)}$ iff $R_{t_2(i,j)}(x,t_2(i,j),t_1(i,j))$ iff $M_2(x,i,j,d(t_1(i,j),t_2(i,j)),d(t_2(i,j),t_1(i,j))$ iff $M_2(x,i,j,\phi_1(i,j),\phi_2(i,j)$.
Thus $\omega_{\phi_1(i,j)}=\{x:M_1(x,i,j,\phi_1(i,j),\phi_2(i,j))\}$ and $\omega_{\phi_2(i,j)}=\{x:M_2(x,i,j,\phi_1(i,j),\phi_2(i,j))\}$, thus completing the proof.

\section*{Problem 3}
(a): Let $R_1(x,y,y_1,z_1,z_2)$ be $M_1(x,y,z_1,z_2)$, and $R_2(x,y,y_1,z_1,z_2)$ be $M_2(x,y_1,z_1,z_2)$. By Theorem 2, there are recursive functions
$\omega_{\phi_1(i,j)}$ and $\omega_{\phi_2(i,j)}$ such that:\newline
(1) $\omega_{\phi_1(i,j)}=\{x:R_1(x,i,j,\phi_1(i,j),\phi_2(i,j)\}=\{x:M_1(x,i,\phi_1(i,j),\phi_2(i,j)\}$;\newline
(2) $\omega_{\phi_2(i,j)}=\{x:R_2(x,i,j,\phi_1(i,j),\phi_2(i,j)\}=\{x:M_2(x,j,\phi_1(i,j),\phi_2(i,j)\}$\newline thus completing the proof of Theorem 3.


(b): Let $M(x,y,z_1,z_2)$ be the relation $R_y(x,z_1,z_2)$ (as a relation between $x,y,z_1,z_2)$. 
Thus, by Theorem 3, there are recursive functions $\phi_1(i,j)$ and $\phi_2(i,j)$ such that:\newline
(1) $\omega_{\phi_1(i,j)}=\{x:R_i(x,\phi_1(i,j),\phi_2(i,j))\}$;\newline
(2) $\omega_{\phi_2(i,j)}=\{x:R_j(x,\phi_1(i,j),\phi_2(i,j))\}$\newline thus showing that Theorem 4 is a special case of Theorem 3.

\end{document}
